\section{Checked zero-free widths and the smoothed explicit formula}

This section records the analytic inputs now supplied by the checked
RiemannGaussian extraction.  Its declarations retain their upstream
namespaces; their project modules lie under
\texttt{Sidon/NumberTheory/RiemannGaussian}.
The exact source revision, license, modifications and declaration
dependencies are recorded in the extraction manifest.
All statements below concern mathlib's actual zeta function and its
analytic multiplicities.

\begin{definition}[The actual zero divisor]
\label{analytic:zero-model}
\lean{RiemannGaussian.IsNontrivialZetaZero,RiemannGaussian.NontrivialZetaZero,RiemannGaussian.analyticZetaZeroMultiplicity}\leanok
A nontrivial zero is a complex number $\rho$ such that
$\zeta(\rho)=0$, $\rho\ne1$, and
$\rho\ne-2(n+1)$ for every natural number $n$.
Its multiplicity is
\[
 m_\rho=\operatorname{analyticOrderNatAt}(\zeta,\rho).
\]
The zero sum below ranges over these distinct zeros and includes both
signs of the imaginary part.
\end{definition}

\begin{lemma}[Geometry, multiplicity and conjugation]
\label{analytic:zero-symmetry}
\uses{analytic:zero-model}
\lean{RiemannGaussian.NontrivialZetaZero.zero_lt_re,RiemannGaussian.NontrivialZetaZero.re_lt_one,RiemannGaussian.NontrivialZetaZero.conjugate_coe,RiemannGaussian.NontrivialZetaZero.conjugate_conjugate,RiemannGaussian.analyticZetaZeroMultiplicity_positive,RiemannGaussian.analyticZetaZeroMultiplicity_conjugate,RiemannGaussian.NontrivialZetaZero.im_ne_zero_of_eta_mass}\leanok
Every nontrivial zero has $0<\Re\rho<1$, nonzero imaginary part, and
positive finite analytic multiplicity.  Literal complex conjugation
preserves the zero type and its multiplicity and is an involution.
\end{lemma}
\begin{proof}
Nonvanishing of $\zeta$ on $\Re s\ge1$ gives $\Re\rho<1$.
Apply the functional equation to the nontrivial zero $1-\rho$ to
obtain $0<\Re\rho$.  Analyticity near $\rho$ and nonvanishing of
zeta as an analytic function make its order finite; its value zero
makes that order positive.

The identity $\zeta(\overline s)=\overline{\zeta(s)}$ constructs
the conjugate zero.  Conjugation twice is the identity.
Conjugating all derivatives preserves the first nonzero derivative,
so the analytic orders and their natural multiplicities agree.
For nonreality, the positive eta-measure representation gives a
strictly positive tilted cosine moment at ordinate zero whenever
the real part is positive.  At an actual zeta zero that same moment
must be zero.  A real nontrivial zero would therefore contradict
strict positivity.  This argument requires no numerical zero table.
\end{proof}

\begin{lemma}[The checked near-one growth estimate]
\label{analytic:near-one-growth}
\lean{RiemannGaussian.ZetaVinogradovSummedBound.prefix_bound,RiemannGaussian.ZetaVinogradovSummedBound.bound_strip,RiemannGaussian.ZetaVinogradovSummedBound.bound_strip_abs}\leanok
Put
\[
 \Delta_n=\frac1{4096(2n+1)^2},\qquad
 g_n=40\Delta_n^{3/2}.
\]
For $n\ge48$, on the checked line
$1-\Delta_n\le\Re s\le3/2$, above the finite height
$H(8n)$ specified by the degree schedule, one has
\[
 |\zeta(s)|\le
 1048576\,|\Im s|^{g_n}
                 \bigl(\log|\Im s|\bigr)^{2/3}.
\]
\end{lemma}
\begin{proof}
Choose the canonical power-of-two Dirichlet cutoff
$M=2^{\operatorname{depth}(s)+1}$.
Its prefix is the sum of its original dyadic blocks.
The checked finite exponential-sum estimate, with the canonical
scale feasibility inequalities, bounds block $j$ by
\[
 512\,t^{g_n}
 \exp\left(-\frac{(j\log2)^3}{2097152(\log t)^2}\right),
 \qquad t=\Im s>0.
\]
All dyadic endpoints and original coefficients are retained.
The nonnegative cubic-decay envelope has finite partial sums at most
$1024(\log t)^{2/3}$.  The triangle inequality therefore bounds the
entire prefix by
$524288\,t^{g_n}(\log t)^{2/3}$.

The direct Euler continuation identity reconstructs zeta from this
prefix, its complex pole endpoint and its full oscillatory remainder.
At the canonical cutoff their combined norm is at most six.
The height schedule implies $t\ge16$; hence $t^{g_n}\ge1$ and
$(\log t)^{2/3}\ge1/2$.  Doubling the prefix coefficient pays six.
Conjugation handles negative ordinates without altering the real
part, norm or absolute height.
The finite exponential-sum bound and the Euler remainder are proved
declarations in the selected closure, rather than hypotheses of
this statement.
\end{proof}

\begin{theorem}[Vinogradov--Korobov strip]
\label{analytic:vk-strip}
\uses{analytic:zero-symmetry,analytic:near-one-growth}
\lean{RiemannGaussian.ZetaVinogradovSummedCost.margin_of_budget,RiemannGaussian.ZetaVinogradovSummedCost.cost_tendsto,RiemannGaussian.ZetaVinogradovSummedZeroFree.exists_eventual_strip}\leanok
For every $0<C<3\pi/10640$, there is a finite $U\ge4$ such that
every nontrivial zero $\rho$ with $|\Im\rho|\ge U$ satisfies
\[
 w(C,\Im\rho)<\Re\rho<1-w(C,\Im\rho),\qquad
 w(C,t)=
 \frac{C}{(\log(|t|+2))^{2/3}
             (\log\log(|t|+2))^{1/3}}.
\]
\end{theorem}
\begin{proof}
Write $L(t)=\log(|t|+2)$ and $\ell(t)=\log L(t)$.
The checked cubic schedule chooses a natural degree $n(t)$ tending
to infinity and proves its height condition eventually.
At this schedule the normalized radius and growth estimates give
\[
 \frac{w(C,t)\ell(t)}{\Delta_{n(t)}}\longrightarrow4C,
 \qquad
 \frac{g_{n(t)}L(t)}{\ell(t)}\longrightarrow5.
\]
The near-one logarithmic profile is
\[
 P_n(t)=\log1048576+g_nL(t)+\tfrac23\ell(t).
\]
Consequently $P_{n(t)}(t)/\ell(t)$ and
$P_{n(t)}(2t)/\ell(t)$ both tend to $17/3$.
The center shift is $x=6w(C,t)$.
Since $\log(1+1/x)/\ell(t)\to2/3$, the two normalized allowances
$P_n(t)+\log(1+1/x)$ and
$P_n(2t)+\log(1+1/x)$ tend to $19/3$.

The actual signed local detector uses
\[
 B(n,x,t)=1344\,\operatorname{localZetaLogHeight}(0)
 +\frac{8(P_n(t)+\log(1+1/x))
          +2(P_n(2t)+\log(1+1/x))}
        {\pi\Delta_n}.
\]
The preceding limits imply
\[
 w(C,t)B(n(t),6w(C,t),t)\longrightarrow
 \frac{760C}{3\pi}.
\]
Also $w(C,t)\to0$ and
$w(C,t)/\Delta_{n(t)}\to0$.  Thus the complete payment
\[
 Q(C,t)=14w(C,t)B(n(t),6w(C,t),t)
          +392\frac{w(C,t)^2}{\Delta_{n(t)}^2}
\]
tends to $10640C/(3\pi)<1$.

Choose one starting threshold beyond which all the following hold:
$n(t)\ge48$, the checked height condition, $w(C,t)>0$,
$w(C,t)<\Delta_{n(t)}/28$, and $Q(C,t)<1$.
The local detector, applied with its proved analytic-disc and
magnitude bounds at $t$ and $2t$, then gives
$w(C,t)<1-\Re\rho$ for every eligible zero.
All the budget inequalities have been discharged by these limits.

The map $\rho\mapsto1-\overline\rho$ preserves the imaginary part
and replaces the real part by $1-\Re\rho$.
Applying the same right-margin inequality to this zero supplies
the left margin.  The width and payment are invariant under
changing the sign of the ordinate, so the assertion holds at both
signs above the same finite threshold.
\end{proof}

\begin{theorem}[One width for every zero up to height $T$]
\label{analytic:uniform-zero-width}
\uses{analytic:vk-strip}
\lean{Sidon.logarithmicZeroFreeWidth_pos,Sidon.logarithmicZeroFreeWidth_antitone,Sidon.tendsto_logarithmicZeroFreeWidth_zero,Sidon.logarithmicZeroFreeWidth_shift_le,Sidon.exists_logarithmicZeroFreeWidth_of_height_margin,Sidon.exists_logarithmicZetaZeroFreeWidth}\leanok
There is $T_0$ such that for every $T\ge T_0$ and every actual
positive-height zero occurrence $\rho$ up to height $T$,
\[
 \Re\rho\le1-\eta(T),\qquad
 \eta(T)=
 \frac{1/20000}{(\log T)^{2/3}(\log\log T)^{1/3}}.
\]
Occurrences retain the analytic multiplicity.
\end{theorem}
\begin{proof}
Take $C_0=1/10000<3\pi/10640$ and enlarge the preceding threshold
to $U\ge e^2$.
On positive ordinates its source width equals the unshifted project
width with coefficient $C_0$ at $\gamma+2$.
Both logarithmic factors are positive and increasing in this range,
so for $U\le\gamma\le T$ the width at $\gamma+2$ is at least
the width at $T+2$.

For $T\ge2$ and $\log T\ge2$, one has $T+2\le T^2$ and hence
\[
 \log(T+2)\le2\log T,\qquad
 \log\log(T+2)\le\log2+\log\log T\le2\log\log T.
\]
The denominator increases by at most
$2^{2/3}2^{1/3}=2$; the shifted source width therefore dominates
$\eta(T)$.

It remains to cover $0<\gamma\le U$.
The closed discrete zero set meets a compact disk in finitely many
points, so there are finitely many such critical-strip zeros.
Each has $\Re\rho<1$.
Since $\eta(T)\to0$, it is eventually smaller than each of these
finitely many positive margins $1-\Re\rho$.
Taking the maximum of these thresholds proves the assertion.
Repeating a point by its multiplicity preserves its real-part bound.
\end{proof}

\begin{lemma}[Complete reciprocal-square zero mass]
\label{analytic:inverse-square-mass}
\uses{analytic:zero-symmetry}
\lean{RiemannGaussian.RosserSchoenfeldZeroMass.hasSum_poisson,RiemannGaussian.RosserSchoenfeldZeroMass.poisson_lower,RiemannGaussian.RosserSchoenfeldZeroMass.norm_square_mass_summable}\leanok
The actual nontrivial zero divisor has finite mass
\[
 M=\sum_\rho\frac{m_\rho}{|\rho|^2}<\infty.
\]
\end{lemma}
\begin{proof}
The complete reflected xi logarithmic-derivative expansion makes
the nonnegative Poisson series
\[
 \sum_\rho m_\rho P(\rho),\qquad
 P(\rho)=\frac{1-\beta}{(1-\beta)^2+\gamma^2}
             +\frac{\beta}{\beta^2+\gamma^2},
 \quad \rho=\beta+i\gamma,
\]
summable.  Its multiplicities are the analytic zeta multiplicities.
The numerators are positive and $0<\beta<1$.
Replacing either real-part square in a denominator by one increases
that denominator, so
$P(\rho)\ge(1-\beta)/(1+\gamma^2)+\beta/(1+\gamma^2)
=1/(1+\gamma^2)$.

Closedness of the zeta-zero set and $\zeta(0)=-1/2\ne0$ give
$\epsilon>0$ such that the ball of radius $\epsilon$ about zero
contains no zero.  Thus $|\rho|\ge\epsilon$.
With $D=(1+\epsilon^2)/\epsilon^2$, the inequalities
$\gamma^2\le|\rho|^2$ and $\epsilon^2\le|\rho|^2$ imply
$1+\gamma^2\le D|\rho|^2$.
It follows that
\[
 \frac{m_\rho}{|\rho|^2}
 \le D\frac{m_\rho}{1+\gamma^2}\le Dm_\rho P(\rho).
\]
Comparison of nonnegative series proves the conclusion.
No evaluated small-zero certificate or numerical mass constant
is used.
\end{proof}

\begin{definition}[The three logarithmic primitives]
\label{analytic:log-primitives}
\uses{analytic:zero-model}
\lean{RiemannGaussian.RosserSchoenfeldPrimePrimitive.term,RiemannGaussian.RosserSchoenfeldPrimePrimitive.value,RiemannGaussian.RosserSchoenfeldZeroPrimitive.atom,RiemannGaussian.RosserSchoenfeldZeroPrimitive.term,RiemannGaussian.RosserSchoenfeldZeroPrimitive.value,RiemannGaussian.RosserSchoenfeldTrivialPrimitive.value}\leanok
For $t\ge0$ put
\[
 A(t)=\sum_{1\le n\le\lfloor e^t\rfloor}
           \Lambda(n)\max(t-\log n,0),
\quad
 Z(t)=\sum_\rho m_\rho\frac{e^{\rho t}-1}{\rho^2},
\quad
 V(t)=\sum_{j\ge1}\frac{1-e^{-2jt}}{4j^2}.
\]
Their one-sided Laplace transform is
$\mathcal Lf(s)=\int_0^\infty e^{-st}f(t)\,dt$.
\end{definition}

\begin{lemma}[Continuity, convergence and exact transforms]
\label{analytic:primitive-transforms}
\uses{analytic:log-primitives,analytic:inverse-square-mass}
\lean{RiemannGaussian.RosserSchoenfeldPrimePrimitive.continuous_value,RiemannGaussian.RosserSchoenfeldPrimePrimitive.transform_value,RiemannGaussian.RosserSchoenfeldZeroPrimitive.summable_norm_term,RiemannGaussian.RosserSchoenfeldZeroPrimitive.integrable_value,RiemannGaussian.RosserSchoenfeldZeroPrimitive.transform_value,RiemannGaussian.RosserSchoenfeldTrivialPrimitive.transform_value}\leanok
All three primitives are continuous at positive time.
Their infinite series converge absolutely; after damping by any
$s$ with $\Re s>1$, all required absolute integrals are summable.
Their exact transforms are
\[
 \mathcal LA(s)=-\frac{\zeta'(s)}{s^2\zeta(s)},\quad
 \mathcal LZ(s)=
 \frac{(\xi'/\xi)(s)-(\xi'/\xi)(0)}{s^2},\quad
 \mathcal LV(s)=
 \frac{\psi(1+s/2)+\gamma_E}{2s^2}.
\]
Here $\psi=\Gamma'/\Gamma$.
\end{lemma}
\begin{proof}
The arithmetic hinge for $n$ vanishes outside its moving finite window.
On a neighborhood of a fixed time, one larger finite window includes
every possible contributing index.  Each hinge is continuous even
where its argument is zero, proving continuity of $A$.
For $\Re s>1$, the sum of the absolute integrals is bounded by a
convergent von Mangoldt Dirichlet series divided by $(\Re s)^2$.
Hence sum and integral may be interchanged.  Substitution
$u=t-\log n$ evaluates one hinge transform as
$\Lambda(n)n^{-s}/s^2$.  The absolutely convergent Euler-product
logarithmic derivative gives the transform of $A$.

For a nontrivial zero, $|e^{\rho t}|\le e^t$, whence
\[
 \left|m_\rho\frac{e^{\rho t}-1}{\rho^2}\right|
 \le(e^t+1)\frac{m_\rho}{|\rho|^2}.
\]
The preceding mass lemma proves absolute and locally uniform
convergence, continuity, and the bound $|Z(t)|\le M(e^t+1)$.
After damping, the majorant is
$M(e^{(1-\Re s)t}+e^{-\Re(s)t})$, which is integrable.
Thus the full series may again be integrated termwise.
Since $\rho\ne0$ and $\Re(s-\rho)>0$,
\[
 \int_0^\infty e^{-st}\frac{e^{\rho t}-1}{\rho^2}\,dt
 =\frac1{\rho^2}\left(\frac1{s-\rho}-\frac1s\right)
 =\frac1{s\rho(s-\rho)}.
\]
Summing the complete checked xi logarithmic-derivative difference
gives the claimed transform of $Z$.

Each trivial-zero term is nonnegative and at most $1/(4j^2)$.
That summable majorant has total $\pi^2/24$ and proves uniform
convergence, continuity including at zero, and integrability after
positive damping.  Its single-term transform is
$1/(2js(s+2j))$.  The convergent digamma difference series sums
these terms to the displayed transform of $V$.
Every exchange uses the stated absolute majorants; no symmetric
cutoff interpretation is introduced.
\end{proof}

\begin{lemma}[Uniqueness from one vertical transform line]
\label{analytic:laplace-uniqueness}
\lean{RiemannGaussian.RosserSchoenfeldLaplace.zero_of_vertical_transform}\leanok
Suppose $f$ is continuous on positive time,
$e^{-at}f(t)$ is integrable there, and
$\mathcal Lf(a+iy)=0$ for every real $y$.
Then $f(t)=0$ for every $t>0$.
\end{lemma}
\begin{proof}
Extend $F(t)=e^{-at}f(t)$ by zero to $t\le0$.
It is integrable on the real line and continuous at every positive
time.  Its Fourier transform at frequency $w$ equals
$\mathcal Lf(a+2\pi iw)$ and therefore vanishes.
The zero transform is integrable.
Fourier inversion at each positive continuity point gives $F(t)=0$.
The exponential factor never vanishes, so $f(t)=0$.
\end{proof}

\begin{theorem}[The literal logarithmically smoothed explicit formula]
\label{analytic:smoothed-explicit-formula}
\uses{analytic:primitive-transforms,analytic:laplace-uniqueness}
\lean{RiemannGaussian.RosserSchoenfeldExplicitFormula.prime_eq_spectral,RiemannGaussian.RosserSchoenfeldExplicitFormula.finite_prime_explicitFormula}\leanok
For every $x\ge1$, in $\mathbb C$,
\[
 \sum_{1\le n\le\lfloor x\rfloor}\Lambda(n)\log(x/n)
 =x-1-\log(2\pi)\log x-Z(\log x)+V(\log x).
\]
The zero primitive includes every nontrivial zero with actual
analytic multiplicity and both ordinate signs.
\end{theorem}
\begin{proof}
Set
$B(t)=e^t-1-\log(2\pi)t-Z(t)+V(t)$.
Its elementary part has transform
$1/(s(s-1))-\log(2\pi)/s^2$.
The exact gamma factorization gives
\[
 \frac{\xi'}{\xi}(s)=\frac{\zeta'}{\zeta}(s)
 +\frac1{s-1}+\frac1s-\frac{\log\pi}{2}
 +\frac12\psi(s/2),
\quad
 \frac{\xi'}{\xi}(0)=
 -1-\frac{\gamma_E}{2}+\log2+\frac{\log\pi}{2}.
\]
Together with $\psi(1+s/2)=\psi(s/2)+2/s$ and the preceding
three transforms, substitution cancels every pole, gamma and
constant term and yields
$\mathcal LB(s)=-(\zeta'/\zeta)(s)/s^2=\mathcal LA(s)$
for $\Re s>1$.
All damped expressions are genuinely integrable.

Apply vertical-line uniqueness to $f=A-B$ on $\Re s=2$.
This proves $A(t)=B(t)$ for $t>0$.
At $t=0$ the arithmetic primitive is zero since $\Lambda(1)=0$;
each zero and trivial-zero term also vanishes, as does $e^t-1$.
Thus equality includes the endpoint.

Finally put $t=\log x$.
For every positive index in the finite window,
$t-\log n\ge0$ and
$t-\log n=\log(x/n)$.
The moving window is exactly $1\le n\le\lfloor x\rfloor$.
Replacing the hinges by these linear weights gives the asserted
formula, including $x=1$.
\end{proof}


\begin{lemma}[A crude multiplicity-sensitive total count]
\label{analytic:crude-zero-count}
\uses{analytic:zero-model}
\lean{Sidon.riemannXi_fourThirdsGrowth,Sidon.jensen_riemannXi_divisor_le_fourThirds,Sidon.card_zetaZeroOccurrences,Sidon.zetaZeroCount_zero_le_riemannXi_divisor,Sidon.zetaZeroCount_zero_le_fourThirds,Sidon.zetaZeroCount_le_fourThirds_add_one}\leanok
There is an absolute constant $C>0$ such that, for $T\ge1$ and
$\sigma\ge0$,
\[
 N(\sigma,T)\le C T^{4/3}.
\]
For every $T\ge0$ the same constant gives
$N(\sigma,T)\le C(T+1)^{4/3}$.
\end{lemma}
\begin{proof}
The theta-tail Mellin estimate already gives constants $C_0\ge1$,
$p>0$ and $B\ge0$ such that, with $R=\lVert z\rVert+1$,
\[
 \lVert\zeta_0^*(z)\rVert
 \le \frac{2C_0}{p}\exp(BR^{4/3}).
\]
Here $\zeta_0^*$ is the entire completed-zeta function used to define
$\xi(z)=z(1-z)\zeta_0^*(z)-1$.
Put $K=\max(1,2C_0/p)$ and $A=K+B+4$.
Since $R\ge1$, both $R\le R^{4/3}$ and $1\le R^{4/3}$ hold.
The elementary inequalities
$2\le e$, $K\le e^K$ and $R^2\le e^{2R}$ show
\[
 \lVert\xi(z)\rVert
 \le K R^2 e^{BR^{4/3}}+1
 \le 2K R^2 e^{BR^{4/3}}
 \le \exp(A R^{4/3}).
\]
In the middle step the product is at least one.
For the last step,
$1+K+2R+BR^{4/3}\le A R^{4/3}$.
Thus the existing Mellin exponent is retained without weakening it.

Apply Jensen's inequality to the entire function $\xi$ on the
concentric disks of radii $r$ and $2r$.
The central value is $\xi(0)=-1$, so its logarithmic contribution
vanishes, and the annular logarithm is $\log2>0$.
The growth bound on the outer circle therefore gives
\[
 \sum_{|z|\le r}\operatorname{ord}_z(\xi)
 \le \frac{A(2r+1)^{4/3}}{\log2}.
\]
The checked sum is the finite sum of the analytic divisor, with no
enumeration or simplicity assumption on zeros.

For each distinct zero in the positive-height critical rectangle,
the occurrence set has exactly one pair $(\rho,k)$ for each integer
$0\le k<\operatorname{ord}_\rho(\zeta)$.
Fibers belonging to different complex zeros are disjoint.
Consequently its cardinality equals the sum of these analytic
multiplicities.
Every such zero is nontrivial: its positive imaginary part excludes
both negative even integers and the pole at one.
The checked completed-zeta factorization identifies its multiplicity
with the analytic order of $\xi$.
Moreover $0\le\beta\le1$ and $0<\gamma\le T$ imply
$\lVert\rho\rVert\le|\beta|+|\gamma|\le T+1$.
Thus the occurrence count is bounded by the divisor sum on the disk
of radius $r=T+2$.

For $T\ge1$, $2(T+2)+1=2T+5\le7T$.
The preceding Jensen estimate is therefore at most
\[
 \frac{A7^{4/3}}{\log2}\,T^{4/3}.
\]
Filtering further by $\beta\ge\sigma\ge0$ only decreases the count.
Finally monotonicity in height bounds $N(\sigma,T)$ by
$N(\sigma,T+1)$ and gives the stated bound for every $T\ge0$.
The exponent $4/3$ is strictly less than $23/16$, leaving a positive
power reserve for the low-real-part range of the density estimate.
\end{proof}


\begin{lemma}[Vanishing of the actual weighted zero sum]
\label{analytic:weighted-zero-vanishing}
\uses{construction:zero-layers,construction:logarithmic-reserve,analytic:uniform-zero-width}
\lean{Sidon.eventually_logarithmic_width_floor_reserve,Sidon.eventually_zetaZeroWeight_le_inverse_log_sq,Sidon.tendsto_zetaZeroWeight_of_density_width,Sidon.tendsto_zetaZeroWeight_rpow_of_density_width,Sidon.tendsto_zetaZeroWeight_seventeen_fiftieths_of_density,Sidon.eventually_zetaZeroWeight_seventeen_fiftieths_le}\leanok
Suppose the actual positive-height zero count satisfies the uniform bound
\[
 N(\sigma,T)\le K T^{A(1-\sigma)}(\log T)^B
 \qquad(0\le\sigma\le1,\ T\ge T_0),
\]
where $A,K\ge0$ and $B$ is a nonnegative integer.
Suppose also that every occurrence up to height $T$ has
$\beta\le1-\eta(T)$ eventually, where
\[
 \eta(T)=\frac{C}{(\log T)^{2/3}(\log\log T)^{1/3}},
 \qquad C>0.
\]
For every fixed $r>A$,
\[
 \sum_{(\rho,k)\in\mathcal O(T)}
       \exp\bigl(-r\log T(1-\Re\rho)\bigr)\longrightarrow0.
\]
The same conclusion holds with $T=x^\tau$ and weight
$\exp(-\log x(1-\Re\rho))$ whenever $\tau>0$ and $A<1/\tau$.
In particular $\tau=17/50$ is admissible for $A=23/8$.
The checked unconditional zero-free strip supplies the second
hypothesis in the final density-only wrappers.
\end{lemma}
\begin{proof}
Put $L=r\log T$, $b=A/r<1$ and $q=1-b>0$.
For sufficiently large $T$, all quantities being logged are positive.
Partition the finite occurrence set into the layers
$j=\lfloor L(1-\beta)\rfloor$.
The density bound, evaluated at
$\max(0,1-(j+1)/L)$, bounds the cardinality of layer $j$ by
\[
 \frac{K}{r^B}L^B\,e^{b(j+1)}.
\]
The clipped real part lies in $[0,1]$.
Every layer element has $\beta$ at least that clipped value;
the clipped-to-zero case follows from $\beta\ge0$.
Moreover
$A\log T(1-\max(0,1-(j+1)/L))\le b(j+1)$.
These facts justify using the uniform density constant at a real
part which depends on $T$ and $j$.

Every occupied layer satisfies $j\ge\lfloor L\eta(T)\rfloor$,
and every weight in that layer is at most $e^{-j}$.
The finite sum is therefore bounded by the convergent geometric
tail
\[
 \frac{(K/r^B)L^B e^b}{1-e^{-q}}\,
       \exp\bigl(-q\lfloor L\eta(T)\rfloor\bigr).
\]
It remains to absorb $L^B$ uniformly despite the integer floor.

The checked logarithmic reserve says that
$\eta(T)\log T/\log\log T\to\infty$.
Apply it with multiplier $qr$ and requested coefficient
$2(B+2)+q$.
For sufficiently large $T$ this gives
\[
 [2(B+2)+q]\log\log T\le qr\eta(T)\log T.
\]
Increase the threshold so that
$\log\log T\ge1$ and $\log\log T\ge\log r$.
Then $\log L=\log r+\log\log T\le2\log\log T$.
Since $\lfloor u\rfloor>u-1$ for the nonnegative floor used in
the checked occurrence layers,
\[
 (B+2)\log L
 \le 2(B+2)\log\log T
 \le qL\eta(T)-q
 \le q\lfloor L\eta(T)\rfloor.
\]
Exponentiation bounds the geometric tail by
\[
 \frac{(K/r^B)e^b}{1-e^{-q}}\frac1{L^2}.
\]
The denominator is positive because $q>0$.
As $L=r\log T\to\infty$, the upper bound tends to zero;
the finite sum is nonnegative, so the squeeze theorem proves
the first limit.

For $T=x^\tau$, positive $\tau$ makes $T\to\infty$.
Choose $r=1/\tau$ and use
$r\log(x^\tau)=r\tau\log x=\log x$ for $x>0$.
Composition proves the second limit.
For $\tau=17/50$ and $A=23/8$ the required strict inequality is
$23/8<50/17$.
Finally the limit is eventually below every prescribed
$\varepsilon>0$, including the detector's choice $1/16$.
\end{proof}

The kernel dependency audits of the exact formula, complete zero mass,
literal conjugation, nonreality, and VK strip contain only
\texttt{propext}, \texttt{Classical.choice}, and \texttt{Quot.sound}.
The formula does not take an explicit-formula identification or an
unproved prime-error estimate as an argument.
